% Generated draft: requires mathematical and source review.
\documentclass[11pt]{article}
\usepackage[margin=25mm]{geometry}
\usepackage{fontspec}
\setmainfont{Latin Modern Roman}
\usepackage{amsmath,amssymb,mathtools,unicode-math}
\setmathfont{Latin Modern Math}
\usepackage{enumitem,needspace,xurl,hyperref}
\setlength{\parindent}{0pt}
\setlength{\parskip}{5pt}
\setlength{\emergencystretch}{3em}
\newcommand{\N}{\mathbb{N}}
\newcommand{\Z}{\mathbb{Z}}
\newcommand{\Q}{\mathbb{Q}}
\newcommand{\R}{\mathbb{R}}
\newcommand{\C}{\mathbb{C}}
\newcommand{\F}{\mathbb{F}}
\newcommand{\A}{\mathbb{A}}
\newcommand{\PP}{\mathbb{P}}
\newcommand{\E}{\mathbb{E}}
\newcommand{\eps}{\varepsilon}
\newcommand{\dd}{\,\mathrm{d}}
\begin{document}
\section*{UnsolvedMath ID: 30003928}
\textbf{Source problem code:} catalogue entry
\subsection{Definitions and mathematical statement}
Let \(M\) be a complex manifold and let \(w_0\in M\). The source uses the function class \(\operatorname{PSH}_{cb}(M)\), the sets \(A_{cb}(w_0)\) and \(A_{cb}^{e}(w_0)\), and the manifold core \(c(M)\), but supplies no definitions of these objects. In particular, the meanings of the subscripts \(cb\), the superscript \(e\), and the precise convention for the core remain unresolved. They must be recovered from the cited sources before this becomes a fully specified mathematical statement.

With these symbols retained in their intended source meanings, the record asks the following three questions separately:
\begin{enumerate}
\item For every complex manifold \(M\) and every \(w_0\in M\), does
\[
A_{cb}(w_0)=A_{cb}^{e}(w_0)
\]
hold?
\item For every complex manifold \(M\) and every \(w_0\in M\), does
\[
A_{cb}(w_0)\subseteq c(M)
\]
hold?
\item For every complex manifold \(M\), does there exist \(u\in\operatorname{PSH}_{cb}(M)\) such that
\[
u(w)\geq 0\quad\text{for every }w\in M,
\qquad
\{w\in M:u(w)=0\}=c(M)?
\]
Here the displayed set is the zero set of \(u\).
\end{enumerate}
The supplied wording imposes no further restriction on \(M\) or \(w_0\). Whether additional hypotheses or a restricted choice of \(w_0\) are intended cannot be determined from the record. The original inclusion symbol \(\subset\) is rendered as \(\subseteq\) in the supplied cleaned statement; strict inclusion is not asserted here.
\subsection{Short English statement}
The source asks whether two sets associated with bounded plurisubharmonic functions coincide, whether the first lies in the manifold core, and whether that core is exactly the zero set of one nonnegative function in the indicated class. Definitions of the class, the two sets, and the core are missing, so the formulation remains conditional on the intended source conventions.
\subsection{Sources}
\begin{itemize}
\item Local UnsolvedMath registry: supplied record with permanent ID \(30003928\), title \emph{Bounded Plurisubharmonic Functions and Manifold Cores}, source problem code \texttt{catalogue entry}, in the \emph{Oberwolfach Reports Open Problems} collection. The proposed year is recorded as 2018; no individual proposer is supplied.
\item Supplied source citation: \emph{Geometric Methods of Complex Analysis (2019)}, Oberwolfach Reports, pp.\ 2253--2302, DOI \texttt{10.4171/owr/2018/37}. \url{https://doi.org/10.4171/owr/2018/37}. The literature-source list also labels this reference \emph{Oberwolfach Report 37/2018}.
\item Supplied literature reference: \emph{Plurisubharmonically separable complex manifolds}, attributed in the supplied assessment to Poletsky--Shcherbina. \url{https://doi.org/10.1090/proc/14222}.
\item Literature assessment dated exactly \texttt{2026-08-21}: the supplied curation labels the problem open, reports structural results for \(A_{cb}^{e}\) and the core in the Poletsky--Shcherbina paper, and states that the three questions remain open and that no later resolution was located. This is supplied dated curation derived from the OWR record, not an independently verified literature assessment. The record expressly requests independent verification of time-sensitive status and cited claims.
\end{itemize}
\end{document}
